%% Pakiety. Dokument kompilowany przez LuaLaTeX (patrz .latexmkrc).

% Język i typografia
\usepackage[polish]{babel}
\usepackage{fontspec}             % Czcionki OpenType (domyślnie Latin Modern)
\usepackage{microtype}            % Lepsze łamanie i justowanie
\usepackage{csquotes}             % \enquote{...} z cudzysłowami zależnymi od języka

% Układ strony
\usepackage[margin=2.5cm]{geometry}
\usepackage{setspace}             % \onehalfspacing, \doublespacing

% Matematyka i jednostki
\usepackage{amsmath}
\usepackage{amssymb}
\usepackage[locale=PL]{siunitx}   % \qty{3.3}{\volt}, \num{1e-3}

% Rysunki i tabele
\usepackage{graphicx}
\usepackage{float}                % Umiejscowienie [H]
\usepackage{booktabs}             % \toprule, \midrule, \bottomrule
\usepackage{tabularx}
\usepackage[font=small,labelfont=bf]{caption}
\usepackage{subcaption}

% Listy i drobiazgi
\usepackage{enumitem}
\usepackage{xspace}               % Odstęp po makrach tekstowych

% Kolory i grafika wektorowa
\usepackage[table]{xcolor}
\usepackage{tikz}

% Bibliografia (biber)
\usepackage[backend=biber,style=numeric,sorting=none]{biblatex}

% Odnośniki — hyperref ładujemy jako jeden z ostatnich pakietów
\usepackage[hidelinks]{hyperref}
\usepackage[nameinlink]{cleveref} % \cref{...} → „rys. 1”, „tab. 2”
